% GENERATED FILE. Edit canonical lessons, then run npm run generate:books.
% canonical-source-hash: fnv1a64:79938281e299134b
% canonical-lessons: FR-C213-soulever, FR-C213-surprendre, FR-C213-soupconner, FR-C213-serrer, FR-C213-suggerer, FR-R213-first-pass-words-from-chapters-205-208, FR-R213-second-pass-words-from-chapters-209-213

\chapter{To lift, To surprise, To suspect, To squeeze, To suggest}
\label{ch:fr-to-lift-to-surprise-to-suspect-to-squeeze-to-suggest}
\hlchaptermodality{213}

\begin{chapteropening}
\textbf{By the end of this chapter:} \emph{I can say to lift, to surprise, to suspect, to squeeze and to suggest.}

Complete the last lesson of chapter 213: I can say to lift, to surprise, to suspect, to squeeze and to suggest.
\end{chapteropening}

\section[soulever]{soulever — to lift}

\label{lesson:FR-C213-soulever}



\begin{quote}
Before the new one: say the French for to float, then the French for to chase.
\end{quote}

\subsection*{You'll want to know: soulever}
\textbf{soulever} — \textquotedblleft{}to lift\textquotedblright{}.

\begin{grammarlens}[title={What you've built}]
One more word for this chapter.
\end{grammarlens}

\subsection*{Your turn}
\begin{itemize}
\raggedright
  \item \emph{Say it:} \emph{soulever}
  \item \emph{Say it:} \emph{soulever}, once more
  \item \emph{Say it:} say \emph{soulever}
  \item \emph{Recall:} say the French for to float, then the French for to chase, then say \emph{soulever} again
\end{itemize}

\begin{tcolorbox}[breakable,colback=teal!4,colframe=teal!35!black,title={Before you move on}]
What does soulever mean? (\textquotedblleft{}To lift\textquotedblright{}.) Say it once more.
\end{tcolorbox}
\section[surprendre]{surprendre — to surprise}

\label{lesson:FR-C213-surprendre}



\begin{quote}
Before the new one: say the French for to chase, then the French for to lift.
\end{quote}

\subsection*{You'll want to know: surprendre}
\textbf{surprendre} — \textquotedblleft{}to surprise\textquotedblright{}.

\begin{grammarlens}[title={What you've built}]
One more word for this chapter.
\end{grammarlens}

\subsection*{Your turn}
\begin{itemize}
\raggedright
  \item \emph{Say it:} \emph{surprendre}
  \item \emph{Say it:} \emph{surprendre}, once more
  \item \emph{Say it:} say \emph{surprendre}
  \item \emph{Recall:} say the French for to chase, then the French for to lift, then say \emph{surprendre} again
\end{itemize}

\begin{tcolorbox}[breakable,colback=teal!4,colframe=teal!35!black,title={Before you move on}]
What does surprendre mean? (\textquotedblleft{}To surprise\textquotedblright{}.) Say it once more.
\end{tcolorbox}
\section[soupçonner]{soupçonner — to suspect}

\label{lesson:FR-C213-soupconner}



\begin{quote}
Before the new one: say the French for to lift, then the French for to surprise.
\end{quote}

\subsection*{You'll want to know: soupçonner}
\textbf{soupçonner} — \textquotedblleft{}to suspect\textquotedblright{}.

\begin{grammarlens}[title={What you've built}]
One more word for this chapter.
\end{grammarlens}

\subsection*{Your turn}
\begin{itemize}
\raggedright
  \item \emph{Say it:} \emph{soupçonner}
  \item \emph{Say it:} \emph{soupçonner}, once more
  \item \emph{Say it:} say \emph{soupçonner}
  \item \emph{Recall:} say the French for to lift, then the French for to surprise, then say \emph{soupçonner} again
\end{itemize}

\begin{tcolorbox}[breakable,colback=teal!4,colframe=teal!35!black,title={Before you move on}]
What does soupçonner mean? (\textquotedblleft{}To suspect\textquotedblright{}.) Say it once more.
\end{tcolorbox}
\section[serrer]{serrer — to squeeze}

\label{lesson:FR-C213-serrer}



\begin{quote}
Before the new one: say the French for to surprise, then the French for to suspect.
\end{quote}

\subsection*{You'll want to know: serrer}
\textbf{serrer} — \textquotedblleft{}to squeeze\textquotedblright{}.

\begin{grammarlens}[title={What you've built}]
One more word for this chapter.
\end{grammarlens}

\subsection*{Your turn}
\begin{itemize}
\raggedright
  \item \emph{Say it:} \emph{serrer}
  \item \emph{Say it:} \emph{serrer}, once more
  \item \emph{Say it:} say \emph{serrer}
  \item \emph{Recall:} say the French for to surprise, then the French for to suspect, then say \emph{serrer} again
\end{itemize}

\begin{tcolorbox}[breakable,colback=teal!4,colframe=teal!35!black,title={Before you move on}]
What does serrer mean? (\textquotedblleft{}To squeeze\textquotedblright{}.) Say it once more.
\end{tcolorbox}
\section[suggérer]{suggérer — to suggest}

\label{lesson:FR-C213-suggerer}



\begin{quote}
Before the new one: say the French for to suspect, then the French for to squeeze.
\end{quote}

\subsection*{You'll want to know: suggérer}
\textbf{suggérer} — \textquotedblleft{}to suggest\textquotedblright{}.

\begin{grammarlens}[title={What you've built}]
One more word for this chapter.
\end{grammarlens}

\subsection*{Your turn}
\begin{itemize}
\raggedright
  \item \emph{Say it:} \emph{suggérer}
  \item \emph{Say it:} \emph{suggérer}, once more
  \item \emph{Say it:} say \emph{suggérer}
  \item \emph{Recall:} say the French for to suspect, then the French for to squeeze, then say \emph{suggérer} again
\end{itemize}

\begin{tcolorbox}[breakable,colback=teal!4,colframe=teal!35!black,title={Before you move on}]
What does suggérer mean? (\textquotedblleft{}To suggest\textquotedblright{}.) Say it once more.
\end{tcolorbox}
\section[(dialogue)]{Words from chapters 205-208 — a first pass}

\label{lesson:FR-R213-first-pass-words-from-chapters-205-208}



\begin{quote}
Say the words for to squeeze and to suggest.
\end{quote}

\subsection*{Your turn}
Say these aloud:

\begin{itemize}
\raggedright
  \item to book, to spend, to save, to order, to dry
  \item to break, to touch, to taste, to heat, to cool down
  \item to be born, to die, to grow old, to marry, to divorce
  \item to applaud, to cross, to increase, to decrease, to warn
\end{itemize}

\begin{tcolorbox}[breakable,colback=teal!4,colframe=teal!35!black,title={Before you move on}]
To book? (\textbf{réserver}.) To spend? (\textbf{dépenser}.) To save? (\textbf{économiser}.) To order? (\textbf{commander}.) To dry? (\textbf{sécher}.) To break? (\textbf{casser}.) To touch? (\textbf{toucher}.) To taste? (\textbf{goûter}.) To heat? (\textbf{chauffer}.) To cool down? (\textbf{refroidir}.) To be born? (\textbf{naître}.) To die? (\textbf{mourir}.) To grow old? (\textbf{vieillir}.) To marry? (\textbf{épouser}.) To divorce? (\textbf{divorcer}.) To applaud? (\textbf{applaudir}.) To cross? (\textbf{traverser}.) To increase? (\textbf{augmenter}.) To decrease? (\textbf{diminuer}.) To warn? (\textbf{prévenir}.)
\end{tcolorbox}
\section[(dialogue)]{Words from chapters 209-213 — a second pass}

\label{lesson:FR-R213-second-pass-words-from-chapters-209-213}



\begin{quote}
Say the words for to enjoy and to imagine.
\end{quote}

\subsection*{Your turn}
Say these aloud:

\begin{itemize}
\raggedright
  \item to enjoy, to imagine, to point out, to inform, to complain
  \item to reach, to achieve, to give a present, to give back, to recognise
  \item to be silent, to fear, to attempt, to transport, to kill
  \item to fix, to brake, to function, to float, to chase
  \item to lift, to surprise, to suspect, to squeeze, to suggest
\end{itemize}

\begin{tcolorbox}[breakable,colback=teal!4,colframe=teal!35!black,title={Before you move on}]
To enjoy? (\textbf{profiter}.) To imagine? (\textbf{imaginer}.) To point out? (\textbf{indiquer}.) To inform? (\textbf{informer}.) To complain? (\textbf{se plaindre}.) To reach? (\textbf{atteindre}.) To achieve? (\textbf{réaliser}.) To give a present? (\textbf{offrir un cadeau}.) To give back? (\textbf{rendre}.) To recognise? (\textbf{reconnaître}.) To be silent? (\textbf{se taire}.) To fear? (\textbf{craindre}.) To attempt? (\textbf{tenter}.) To transport? (\textbf{transporter}.) To kill? (\textbf{tuer}.) To fix? (\textbf{fixer}.) To brake? (\textbf{freiner}.) To function? (\textbf{fonctionner}.) To float? (\textbf{flotter}.) To chase? (\textbf{poursuivre}.) To lift? (\textbf{soulever}.) To surprise? (\textbf{surprendre}.) To suspect? (\textbf{soupçonner}.) To squeeze? (\textbf{serrer}.) To suggest? (\textbf{suggérer}.)
\end{tcolorbox}
